%%%%%%%%%%  scan idx 208  |  folio 195  %%%%%%%%%%
\section*{APPENDICE I\\[3pt] UN CONTRE-EXEMPLE DE VERDIER\protect\footnote{Cet appendice reproduit, à quelques détails de rédaction près, deux lettres de J.-L. Verdier (30/9/67 et 11/10/67) au rédacteur, en réponse à une question de celui-ci.}}

\subsection*{0. Introduction}
Si $S$ est un schéma, on note $\Mod(S)$ (resp. $\Qcoh(S)$) la catégorie des $\mathcal O_S$-Modules (resp. $\mathcal O_S$-Modules quasi-cohérents). L'objet de cet appendice est d'attirer l'attention du lecteur sur certaines propriétés pathologiques de la catégorie $\mathrm D(\Qcoh(S))$, et de prouver notamment les faits suivants :

\textbf{0.1.} Soient $X$ un schéma affine d'anneau $A$, et $I$ un $A$-module injectif. On sait ([H] II 7.14) que, si $A$ est noethérien, le faisceau $\widetilde I$ est un $\mathcal O_X$-Module injectif. Si $A$ n'est pas noethérien, $\widetilde I$ n'est pas nécessairement injectif, ni même flasque, même si l'espace sous-jacent à $X$ est noethérien et de dimension de Krull finie. De plus, la restriction de $\widetilde I$ à un ouvert $U$ de $X$ n'est pas en général un injectif de $\Qcoh(U)$.

\textbf{0.2.} Soit $f:X\longrightarrow Y$ un morphisme de schémas quasi-compacts et quasi-séparés. Le foncteur $\mathbb Rf_*:\mathrm D^+(X)\longrightarrow\mathrm D^+(Y)$ n'induit pas en général sur $\mathrm D^+(\Qcoh(X))$ (sous-catégorie pleine de $\mathrm D^+(X)$ d'après (II 3.5)) le dérivé droit de $f_*:\Qcoh(X)\longrightarrow\Mod(Y)$. Plus précisément, un injectif de $\Qcoh(X)$ n'est pas nécessairement acyclique pour $f_*:\Mod(X)\longrightarrow\Mod(Y)$.

\textbf{0.3.} Soit $S$ un schéma quasi-compact et quasi-séparé. Un objet de $\Qcoh(S)$, même injectif, n'est pas nécessairement acyclique pour le cohéreur (II 3.2). Le foncteur $\varphi:\mathrm D^b(\Qcoh(S))\longrightarrow\mathrm D^b(S)$ (II 3.0) n'est pas en général pleinement fidèle.

\subsection*{1. Les contre-exemples}
Rappelons d'abord le lemme suivant (SGA 2 II 9) :

\textbf{Lemme 1.1.} Soient $A$ un anneau, $X=\Spec(A)$, $\mathbf f=(f_\alpha)$ un système fini d'éléments de $A$, $Y$ le fermé défini par $\mathbf f(x)=0$. Pour $i>0$, les conditions suivantes sont équivalentes :
\begin{enumerate}[label=(\roman*)]
\item pour tout $A$-module injectif $I$, on a $H^i_Y(X,\widetilde I)=0$;
\item le pro-objet $n\mapsto H_i(\mathbf f^n,A)$ est nul, i.e. pour tout $n$, il existe $n'\geq n$ 
